% !TeX program = xelatex
% Public student research CV. Compile with XeLaTeX.
\documentclass[11pt,a4paper]{article}
\usepackage[a4paper,left=20mm,right=20mm,top=15mm,bottom=17mm]{geometry}
\usepackage{fix-cm}
\usepackage{fontspec}
\usepackage{xeCJK}
\usepackage{enumitem}
\usepackage{xcolor}
\usepackage{hyperref}
\usepackage{fancyhdr}
\usepackage{microtype}
\usepackage{tabularx}
\usepackage{needspace}

% Explicit TeX Live fonts; no proprietary fonts or synthetic bold.
\setmainfont{texgyretermes-regular.otf}[
  BoldFont=texgyretermes-bold.otf,ItalicFont=texgyretermes-italic.otf,
  BoldItalicFont=texgyretermes-bolditalic.otf]
\setsansfont{texgyretermes-regular.otf}[
  BoldFont=texgyretermes-bold.otf,ItalicFont=texgyretermes-italic.otf,
  BoldItalicFont=texgyretermes-bolditalic.otf]
\setCJKmainfont{FandolSong-Regular.otf}[
  BoldFont=FandolSong-Bold.otf,ItalicFont=FandolKai-Regular.otf]
\setCJKsansfont{FandolHei-Regular.otf}[BoldFont=FandolHei-Bold.otf]

% ===== PERSONAL FIELDS =====
\newcommand{\CVName}{Yiming Zhao}
\newcommand{\CVSecondaryName}{}
\newcommand{\CVUniversity}{The Hong Kong University of Science and Technology (Guangzhou)}
\newcommand{\CVDept}{}
\newcommand{\CVStage}{PhD Student}
\newcommand{\CVResearchLabel}{Research focus:}
\newcommand{\CVResearch}{UAV vision-language navigation, radar-camera fusion, and acoustic sensing.}
\newcommand{\CVEmail}{zym736531683@gmail.com}
\newcommand{\CVAcademicEmail}{yzhao517@connect.hkust-gz.edu.cn}
\newcommand{\CVHomeURL}{https://github.com/zhaoyiming}
\newcommand{\CVCodeURL}{https://github.com/zhaoyiming}
\newcommand{\CVHomeLabel}{Website}
\newcommand{\CVLabel}{Curriculum Vitae}

\definecolor{navy}{HTML}{1A3C5E}
\definecolor{gray3}{HTML}{555555}
\definecolor{gray4}{HTML}{888888}
\definecolor{grayline}{HTML}{B0B8C0}
\hypersetup{colorlinks=true,urlcolor=navy,linkcolor=navy,
  pdfauthor={\CVName},pdftitle={\CVLabel{} --- \CVName}}
\urlstyle{same}
\setlength{\parindent}{0pt}
\setlength{\parskip}{0pt}
\setlength{\emergencystretch}{1.5em}
\hyphenpenalty=3000
\exhyphenpenalty=3000
\raggedbottom

\pagestyle{fancy}
\fancyhf{}
\fancyhead[L]{\rmfamily\footnotesize\color{gray4}\CVName{} --- \CVLabel}
\fancyhead[R]{\rmfamily\footnotesize\color{gray4}\thepage}
\renewcommand{\headrulewidth}{0.3pt}
\setlength{\headheight}{14pt}
\fancypagestyle{firstpage}{\fancyhf{}\renewcommand{\headrulewidth}{0pt}}

% ===== TYPOGRAPHY: edit centrally, never compress individual sections =====
\makeatletter
\renewcommand\normalsize{\@setfontsize\normalsize{11}{14}}
\renewcommand\small{\@setfontsize\small{10}{13}}
\renewcommand\footnotesize{\@setfontsize\footnotesize{9}{11.5}}
\makeatother
\AtBeginDocument{\normalsize}
\newcommand{\cvheadingfont}{\rmfamily}
\newcommand{\cventrygap}{\par\addvspace{6pt}}
\newcommand{\cvsection}[1]{%
  \par\addvspace{10pt}\Needspace{4\baselineskip}%
  {\cvheadingfont\bfseries\fontsize{12}{15}\selectfont\color{navy}#1\par}%
  \nobreak\vspace{3pt}{\color{grayline}\hrule height 0.4pt}%
  \nobreak\vspace{5pt}%
}

% Two-line record: title | location/organization; degree/role | dates.
% The record header stays together; long left fields wrap without hitting dates.
\newcommand{\cvrecord}[4]{%
  \par\Needspace{3\baselineskip}\noindent
  \begin{tabularx}{\textwidth}{@{}>{\raggedright\arraybackslash}X >{\raggedleft\arraybackslash}p{0.32\textwidth}@{}}
  \textbf{#1} & {\small\textcolor{gray3}{#2}}\\
  #3 & {\small\textcolor{gray3}{#4}}
  \end{tabularx}\par
}
\newcommand{\cvaward}[2]{%
  \noindent\begin{tabularx}{\textwidth}{@{}>{\raggedright\arraybackslash}X >{\raggedleft\arraybackslash}p{0.22\textwidth}@{}}
  #1 & {\small\textcolor{gray3}{#2}}
  \end{tabularx}\par
}
\newenvironment{cvitems}{\begin{itemize}[
  leftmargin=1.5em,labelsep=0.55em,itemsep=1pt,topsep=2pt,
  partopsep=0pt,parsep=0pt]}{\end{itemize}}
\newenvironment{cvpublications}{\begin{enumerate}[
  leftmargin=1.6em,labelwidth=1em,labelsep=0.6em,label=\arabic*.,
  itemsep=3pt,topsep=0pt,partopsep=0pt,parsep=0pt,align=left]\raggedright}{\end{enumerate}}
\newenvironment{cvskills}{\par\noindent\tabularx{\textwidth}{@{}>{\bfseries\raggedright\arraybackslash}p{8.8em}>{\raggedright\arraybackslash}X@{}}}{\endtabularx\par}


\newcommand{\cvresearch}[3]{%
  \par\Needspace{5\baselineskip}{\raggedright\textbf{#1}\par}
  \begin{tabularx}{\textwidth}{@{}>{\raggedright\arraybackslash}X >{\raggedleft\arraybackslash}p{0.32\textwidth}@{}}
  #2 & {\small\textcolor{gray3}{#3}}
  \end{tabularx}\par
}
\begin{document}
\thispagestyle{firstpage}
{\centering
  {\cvheadingfont\bfseries\fontsize{24}{28}\selectfont\color{navy}\CVName\par}
  \vspace{3pt}
  {\fontsize{11}{14}\selectfont\CVUniversity\par}
  {\small\CVStage{} \textbar{} Guangzhou, China\par}
  \vspace{2pt}
  {\small\color{gray3}\CVResearchLabel{} \CVResearch\par}
  \vspace{3pt}
  {\small\href{mailto:\CVEmail}{\CVEmail}\enspace\textbar\enspace
    \href{mailto:\CVAcademicEmail}{\CVAcademicEmail}\enspace\textbar\enspace
    \href{\CVCodeURL}{github.com/zhaoyiming}\par}
}

\cvsection{Education}
\cvrecord{The Hong Kong University of Science and Technology (Guangzhou)}{Guangzhou, China}{PhD Student}{Sep 2024 -- present}
{\small IELTS: 6.5; CET-6.\par}
\cventrygap
\cvrecord{Nanjing University of Aeronautics and Astronautics}{Nanjing, China}{Master's in Software Engineering, College of Computer Science}{Sep 2020 -- Apr 2023}
{\small Admitted through postgraduate recommendation; GPA: 3.71/4.00.\par}
\cventrygap
\cvrecord{Henan University}{Kaifeng, China}{Bachelor's in Network Engineering, College of Software}{Sep 2016 -- Jun 2020}
{\small GPA: 3.49/4.00; top 5\%.\par}
\cventrygap
\cvrecord{Yuan Ze University}{Taiwan}{Exchange Student in Information Engineering, College of Information}{Sep 2018 -- Feb 2019}

\cvsection{Publications \& Manuscripts}
% Author order, corresponding-author markers, dates, and submission statuses
% are retained from the supplied September 16, 2026 resume.
% Dates below are source-resume entry dates, not independently verified publication dates.
\begin{cvpublications}
\item \textbf{Yiming Zhao}, Tianshun Li, Jingle He, Ruonan Chai, and Xinhu Zheng*. ``LightVLN: Efficient Aerial Vision-and-Language Navigation with Compact Memory and History-Guided Local Aggregation.'' Sep 2026. Submitted to \textit{IEEE ICRA 2027}.
\item \textbf{Yiming Zhao} and Ying Cui*. ``SQRCF-UAV: Streaming Radar--Camera Fusion for 3D UAV Detection and Estimation.'' Apr 2026. Submitted to \textit{IEEE ICRA 2027}; journal version submitted to \textit{IEEE Transactions on Multimedia}.
\item Zihao Tao, \textbf{Yiming Zhao}, Hongtao Zhao, Zijun Gong, and Ying Cui*. ``Sensing in Low-altitude Wireless Networks: Applications, Systems, Techniques, and Developments.'' Feb 2026. Accepted by \textit{IEEE Communications Magazine}.
\item \textbf{Yiming Zhao}, Ying Cui*, Zijun Gong, and Yang Yang. ``SRCF-UAV: Sparse Radar-Camera Fusion for 3D UAV Detection.'' Aug 2025. Accepted by \textit{IEEE ICRA 2026}; journal version submitted to \textit{IEEE Transactions on Multimedia}.
\item Yanchao Zhao, \textbf{Yiming Zhao}*, Si Li, Hao Han, and Lei Xie. ``UltraSnoop: Placement-agnostic Keystroke Snooping via Ultrasonic Sonar.'' Apr 2023. Published in \textit{ACM Transactions on Internet of Things}.
\item \textbf{Yiming Zhao}, Yanchao Zhao, Huawei Tu*, Qihan Huang, Wenlai Zhao, and Wenhao Jiang. ``Motion Gesture Delimiters for Smartwatch Interaction.'' Feb 2022. Published in \textit{Wireless Communications and Mobile Computing}.
\item N. Zheng, Na Li*, \textbf{Yiming Zhao}, Chao Weng, and Dan Su. ``Tencent Speaker Diarization System for the VoxCeleb Speaker Recognition Challenge 2021.'' Oct 2021. Presented at the \textit{VoxSRC Workshop, Interspeech 2021}.
\end{cvpublications}

\cvsection{Honors \& Awards}
\cvaward{Outstanding Graduate and Advanced Research Individual (postgraduate studies)}{2023}
\cvaward{Tencent Hornbill Elite Talent Program; VoxSRC 2021 Track 4, third place (team)}{2021}
\cvaward{Network Technology Challenge, third place}{2020}
\cvaward{Excellent Student, Scholarship, Outstanding Graduate, and Outstanding Graduation Thesis (undergraduate studies)}{2020}

\newpage
\cvsection{Research Experience}
\cvresearch{LightVLN: Efficient Aerial Vision-and-Language Navigation with Compact Memory and History-Guided Local Aggregation}{UAV vision-language navigation}{Feb 2026 -- present}
\begin{cvitems}
\item Developed a policy with a 0.5B language backbone (1.276B total parameters), single-token historical memory, and history- and instruction-conditioned local aggregation; compressed 256 current-view tokens to 32, with at most 48 observation-derived tokens.
\item Achieved 50.93\%/36.14\% success rates on OpenFly Test-Seen/Test-Unseen and 25.83\% on AerialVLN-S Val-Seen.
\item Validated closed-loop hardware-in-the-loop (HIL) navigation on Jetson Orin NX 16\,GB, reaching 14.61\,Hz model inference and 11.13\,Hz end-to-end decision updates.
\end{cvitems}
\cventrygap
\cvresearch{SRCF-UAV: Sparse Radar-Camera Fusion for 3D UAV Detection}{Multimodal perception and dataset construction}{Jun 2024 -- Sep 2025}
\begin{cvitems}
\item Developed sparse radar-camera fusion using radar/image-assisted query initialization and distance--velocity-aware attention for real-time 3D UAV detection.
\item Built a radar-camera-RTK collection platform and released a centimeter-level multimodal UAV dataset; achieved up to 91.65\% mAP with approximately 17\,ms latency.
\end{cvitems}
\cventrygap
\cvresearch{UltraSnoop: Placement-agnostic Keystroke Snooping via Ultrasonic Sonar}{Acoustic sensing}{Oct 2020 -- Oct 2022}
\begin{cvitems}
\item Used smartphone-generated chirps and keystroke reflections to estimate time differences of arrival, with two virtual coordinate systems and MFCC-based refinement.
\item Achieved placement-agnostic, training-free, and keyboard-independent keystroke inference.
\end{cvitems}

\cvsection{Industry Experience}
\cvrecord{NIO}{Beijing, China}{Speech Algorithms}{Jun 2022 -- May 2024}
\begin{cvitems}
\item Collected speech and source data, generated room impulse responses (RIRs), and constructed an in-car speech enhancement dataset.
\item Improved multichannel speech separation models (RNN-BF and CLDNN) through feature, RIR, and data improvements, increasing wake-up rates by 5\% over existing algorithms.
\item Implemented vehicle-side acoustic echo suppression (AES) and voice activity detection (VAD), including simulation, training, and post-processing; improved TalkNet-based audiovisual lip-movement VAD. Both modules launched on NIO ET5 2.2.
\item Trained and improved Qwen-based LLMs as base decoders for the team's multilingual automatic speech recognition (ASR) system.
\end{cvitems}
\cventrygap
\cvrecord{Tencent AI Lab}{Shenzhen, China}{Speech Algorithm Intern}{May 2021 -- Nov 2021}
\begin{cvitems}
\item Improved voiceprint recognition models (ECAPA-TDNN and X-Vector) using SFace loss, mixture of experts (MoE), and dynamic pruning.
\item Contributed to the team's third-place result in Track 4 of VoxSRC 2021.
\end{cvitems}

\cvsection{Research Skills}
\begin{cvskills}
Speech \& sensing & Acoustic echo suppression, multimodal VAD, speech separation, and time-frequency analysis of speech and radar signals.\\
Computer vision & 2D/3D detection, radar-camera fusion, UAV platforms, RTK-based dataset construction, and TI mmWave radar.\\
LLMs \& navigation & LLM pre-training and supervised fine-tuning, multimodal models, UAV navigation, vision-language reasoning, temporal memory, and action prediction.\\
Frameworks & PyTorch, LoRA, and efficient model training.
\end{cvskills}
\end{document}
